
\section{Thresholds for the variables} \label{ch:thresholds}



Before the compilation of the code (Chapter~\ref{ch:startrun}),
the parameter file ({\tt params.h}) should be consulted to check
whether some vector dimensions are in the desired range.
Most important are
%
\begin{itemize}
\item the maximum particle number {\tt NMAX},
\item the maximum number of regularised KS pairs {\tt KMAX}, and
\item the maximum number of neighbours per particle {\tt LMAX}.
\end{itemize}


The particles are saved in various lists which serve to distinguish
between their funcionality.
The table below explains their nomenclature.
``KS--pairs'' are particles that approach each other in a hyperbolic
encounter; they are given a special treatment by the code (see
Chapter~\ref{ch:regularization}).
If NPAIRS is the amount of KS--pairs, then IFIRST = 2*NPAIRS + 1
is the first single particle (not member of a KS pair), and N the
last one.
NTOT = N + NPAIRS is the total number of particles plus c.m.'s.
Therefore NMAX, the dimension of all vectors containing particle
data should be at least of size N + KMAX, where N is the number of
particles and KMAX the maximum number of expected KS pairs.
If one starts with single particles, KMAX = 10 or 20 should usually
be enough, but in clusters with a large number of primordial binaries,
KMAX must be large.
%
\begin{longtable}{@{}p{1.5cm}p{13.0cm}}
N:      & Total number of particles\\
NBIN0:  & number of primordial binaries (physical bound stars)\\
NBIN:   &  ??? \\
NPAIRS: & Number of binaries (KS--pairs, see Chapter~\ref{ch:regularization}),
          transient unbound pairs as well as persistent binaries\\
NTOT:   & = N + NPAIRS; \\
        & Number of single particles plus centres of masses of
          regularized (KS) pairs \\
KMAX:   & threshold for the amount of allowed KS pairs\\
NMAX:   & = N + KMAX; threshold for the total number of particles and
          the centre of masses
\end{longtable}



\vspace{2cm}
\noindent
\fbox{Hier gibt's noch ein Bildchen!}


